\documentclass[10pt,a4paper]{amsart}
\usepackage[margin=19mm]{geometry}
\usepackage{amsmath,amssymb,mathtools}
\usepackage[scr=rsfs,cal=euler]{mathalfa}
\usepackage{xfrac}
\usepackage[unicode,hidelinks,bookmarks=false]{hyperref}
\newcommand{\ie}{i.e.\ }
\newcommand{\cf}{cf.\ }
\newcommand{\resp}{respectively\ }
\newcommand{\Subsection}{\subsection}
\providecommand{\nobreakdash}{\nobreak\hbox{-}}
\setlength{\emergencystretch}{2em}
\multlinegap0pt
\usepackage{mathtools}
\usepackage{amssymb}
\usepackage{dsfont}
\usepackage{float}
\usepackage{xcolor}
\usepackage{enumerate}
\newtheorem{lemma}{Lemma}
\newtheorem{proposition}{Proposition}
\usepackage{cleveref}
\crefname{lemma}{Lemma}{Lemmas}
\crefname{proposition}{Proposition}{Propositions}
\newcommand{\cC}{\mathcal{C}}
\newcommand{\eps}{\varepsilon}
\DeclareMathOperator{\qSC}{qSC}
\title{Weight distribution bounds to relate minimum distance, list decoding, and symmetric channel performance}
\author{Donald Kougang-Yombi, Jan H{\k{a}}z{\l}a}
\date{Source: arXiv:2604.02994v1 (CC BY 4.0)}
\begin{document}
\maketitle

\noindent\emph{Source and licence.} Weight distribution bounds to relate minimum distance, list decoding, and symmetric channel performance. Version arXiv:2604.02994v1 (CC BY 4.0), distributed by arXiv under the Creative Commons Attribution 4.0 International licence, http://creativecommons.org/licenses/by/4.0/, as recorded in the arXiv licence metadata for this exact version and shown on the abstract page https://arxiv.org/abs/2604.02994v1. Full licence notice: \texttt{licenses/arxiv-2604.02994-CC-BY-4.0.txt}.\par\smallskip

\noindent\textit{Editorial scope.} Counted unit: the article's proposition prop:qsc-vanishing-pb (source0.tex lines 1129--1152). The article's complete original proof of it is delivered with it (source0.tex lines 1154--1214, one \texttt{proof} environment of 61 lines). No mathematics is written, repaired or re-derived: the statement and the proof are the article's own lines, in the article's own notation, and no retained line is substituted or re-flowed.  Retained uncounted as the source-local context the proof needs: lines 1035--1048; lines 1075--1079; lines 1106--1121.  Nothing was omitted from any retained span, so the package carries no omission record.  No retained line contains a citation, so the wrapper carries no bibliography and no bibliography entry.  No retained line contains a figure environment, an image-inclusion command or drawing code, so no image is delivered and the package declares no asset dependency.  Editorial formatting only, in addition: an AMS \texttt{amsart} preamble that restates the article's own declarations and macros used by the retained lines, namely the environments \texttt{lemma}, \texttt{proposition} on separate counters, together with the standard packages those lines need.  The retained environments print in this document as Lemma 1, Proposition 1 in order of appearance, whereas the article prints the same items with the numbering of its own full document.  The complete native source of the article as distributed in the arXiv e-print for this exact version is bundled unchanged as \texttt{sources/197776/source0.tex}.
\begin{lemma}
    \label{lem:qary_balls_intersection}
    Let $q\ge 2$.
    Let $t=p n$ (not necessarily an integer)
    for $0\le p\le 1-1/q$
    and $w=\gamma n$ be an integer
    with $0\le w\le \min(2t,n)$. Then,
\begin{align*}
\mu_q(n,t,w)\le n^2\cdot \exp_q\left(nM^{(q)}(\gamma,p)
\right)\;.
\end{align*}
On the other hand, if $w>2t$, then
$\mu_q(n,t,w)=0$.
\end{lemma}

\begin{lemma}
\label{lem:M_q_continuity}
For every $q\ge 2$, the function $M^{(q)}(\gamma,p)$ is continuous on its domain $\{(\gamma,p): 0\le p\le1-1/q, 0\le \gamma\le\min(2p,1)\}$.
In particular, it is uniformly continuous.
\end{lemma}

\begin{proposition}
\label{prop:poltyrev}
    Let $\cC\subseteq\Sigma^n$ be a $q$-ary
    code with weight
    distribution $(A_w(x))_x$,
    $0<p\le 1-1/q$, $1\le w_0\le n$,
    and $t=p n+n^{\theta}$
    (not necessarily an integer), where $\theta\in(1/2,1)$. Then, 
    for every $x\in\cC$,
    \begin{align}
    \label{eq:22}
        P_{B}(\cC,\qSC_p, x)&\le 2\exp\left(-2n^{2\theta-1}\right)+\sum_{w=1}^{w_0}A_w(x)Z(\qSC_p)^w\\
        &\qquad+\left(\frac{(1-p)(q-1)}{p}\right)^{n^\theta}q^{-nH_q(p)}\sum_{w=w_0+1}^{n}A_w(x)\mu_q(n,t,w)\;.
    \label{eq:23}
    \end{align}
\end{proposition}

\begin{proposition}
\label{prop:qsc-vanishing-pb}
Let $\{\cC_n\}_n$ be a family of $q$-ary codes
with distance $d(\cC_n)\ge\omega(\log(nL))$
for some $L=L(n)$.
Let $\alpha>0$ and $0<p<1-1/q$.

Assume that $A_w(x)\le L(n)$ for every $x\in\cC_n$ and
$w<\alpha n$. Furthermore, assume that
for $\alpha n\le w=\gamma n\le \min(2pn+n^{3/4}, n)$ it holds
\begin{align}
\label{eq:24}
\frac{1}{n}\log_q A_{\gamma n}(x)\le G(\gamma)+o(1)
\end{align}
for some continuous function $G(\gamma)$,
and where $o(1)$ denotes a function that goes to $0$ as
$n$ grows, uniformly for all $w$ and $x\in\cC_n$.

If $F^{(q)}(\gamma, p)>G(\gamma)$ for every
$\alpha\le \gamma\le \min(2p, 1)$, then
\begin{align*}
\lim_{n\to\infty} P_B(\cC_n, \qSC_p)=0\;.
\end{align*}
\end{proposition}

\begin{proof}
Let $t\coloneqq pn+n^{3/4}$ and $x\in\cC_n$. 
Below and for the rest of the proof, 
the $o(1)$ and $o(n)$ terms are uniform in $w$ and $x$.
Applying
\Cref{prop:poltyrev} with $w_0=\lceil \alpha n\rceil-1$, 
it holds
\begin{align*}
P_B(\cC_n,\qSC_p,x)
&\le o(1)+\sum_{d(\cC_n)\le w<\alpha n}A_w(x)Z(\qSC_p)^w
+\exp(o(n))q^{-nH_q(p)}\sum_{w\ge\alpha n}A_w(x)\mu_q(n,t,w)\;.
\end{align*}
We will show that each of the two sums above goes to 0,
uniformly in $x$. First, since $p<1-1/q$, it holds
$Z(\qSC_p)<1$. Furthermore, by assumption, we have
$A_w(x)\le L$ for $w<\alpha n$. Hence,
and using $d(\cC_n)\ge\omega(\log(nL))$,
\begin{align*}
\sum_{d(\cC_n)\le w< \alpha n}
A_w(x)Z(\qSC_p)^w\le nL\cdot Z(\qSC_p)^{d(\cC_n)}=o(1)\;.
\end{align*}
We now turn to the second sum.
Let us write $w=\gamma n$
and $p_n\coloneqq t/n=p+n^{-1/4}$. Using the 
assumption~\eqref{eq:24} and \Cref{lem:qary_balls_intersection},
\begin{align*}
\exp(o(n))q^{-nH_q(p)}\sum_{w\ge\alpha n}A_w(x)\mu_q(n,t,w)&\le
\exp(o(n))\sum_{\alpha n\le w\le \min(2t,n)}
\exp_q\left(n\left[
-H_q(p)+G(\gamma)+M^{(q)}(\gamma,p_n)
\right]\right)\\
&\le
\exp_q\left(o(n)-n\left[
\inf_{\alpha\le \gamma\le\min(2p_n,1)}
H_q(p)-G(\gamma)-M^{(q)}(\gamma,p_n)
\right]\right)\;.
\end{align*}
It remains to show that
$H_q(p)-M^{(q)}(\gamma, p_n)-G(\gamma)$ is positive and uniformly
bounded away from $0$.
For $\alpha\le\gamma\le\min(2p, 1)$ this holds
because, by \Cref{lem:M_q_continuity}, we have
$H_q(p)-M^{(q)}(\gamma,p_n)-G(\gamma)
=F^{(q)}(\gamma,p)-G(\gamma)+o(1)$.
But by assumption $F^{(q)}(\gamma, p)-G(\gamma)>0$
is a strictly positive and continuous function
on the compact interval
$\alpha\le \gamma\le \min(2p,1)$, so it holds
$H_q(p)-M^{(q)}(\gamma,p_n)-G(\gamma)\ge \eps>0$
for $n$ large enough.

For $2p<\gamma\le 2p_n$, 
once again by continuity of $M^{(q)}$
and $G$, we have
$H_q(p)-M^{(q)}(\gamma, p_n)-G(\gamma)
=F^{(q)}(2p,p)-G(2p)+o(1)\ge \eps>0$
for large enough $n$. All in all, we showed
that $P_B(\cC_n,\qSC_p,x)\xrightarrow{n\to\infty}0$
uniformly in $x$, therefore also
$P_B(\cC_n,\qSC_p)\xrightarrow{n\to\infty} 0$.
\end{proof}

\end{document}
